% !TeX program = xelatex
% !TeX encoding = UTF-8
% Standalone document. Compile twice with XeLaTeX (or LuaLaTeX).
% Standard TeX Live / MiKTeX CTeX packages; no external images or bibliography.
% CTeX 自动选择当前系统的中文字体；Windows 使用宋体、黑体与微软雅黑。
\documentclass[UTF8,a4paper,11pt]{ctexart}
\usepackage[left=22mm,right=22mm,top=22mm,bottom=23mm,headheight=15pt]{geometry}
\usepackage{amsmath,booktabs,array,longtable,enumitem,xcolor,fancyhdr}
\usepackage{tikz}
\usetikzlibrary{arrows.meta,positioning,shapes.geometric}
\usepackage[unicode,colorlinks=true,linkcolor=accent,urlcolor=accent]{hyperref}
\definecolor{accent}{HTML}{24566C}
\definecolor{muted}{HTML}{5C6470}
\definecolor{light}{HTML}{F1F6F8}
\hypersetup{pdftitle={第三章学习笔记：知识库构建与智能体化 RAG},pdfauthor={根据李博杰教材整理}}
\ctexset{
  section={format=\Large\sffamily\bfseries\color{accent},beforeskip=18pt,afterskip=8pt},
  subsection={format=\large\sffamily\bfseries,beforeskip=12pt,afterskip=5pt}
}
\setlength{\parindent}{0pt}
\setlength{\parskip}{5pt}
\linespread{1.12}
\setlist{nosep,leftmargin=2em,topsep=4pt,parsep=2pt,itemsep=3pt}
\setlength{\tabcolsep}{6pt}
\renewcommand{\arraystretch}{1.35}
\setlength{\LTpre}{5pt}
\setlength{\LTpost}{8pt}
\newlength{\tableusable}
\newcommand{\code}[1]{{\small\ttfamily #1}}
\pagestyle{fancy}
\fancyhf{}
\fancyhead[L]{\small\color{muted}第三章学习笔记}
\fancyhead[R]{\small\color{muted}知识库与智能体化 RAG}
\fancyfoot[C]{\small\color{muted}\thepage}
\renewcommand{\headrulewidth}{0.3pt}
\setlength{\emergencystretch}{2em}
\widowpenalty=10000
\clubpenalty=10000
\displaywidowpenalty=10000
\begin{document}
\begin{center}
  {\LARGE\sffamily\bfseries\color{accent}第三章学习笔记\par}
  \vspace{5pt}
  {\Large\sffamily\bfseries 知识库构建与智能体化 RAG\par}
\end{center}
\vspace{2pt}
\noindent\textcolor{accent}{\rule{\textwidth}{0.7pt}}


\begin{quote}\small\color{muted}

来源：李博杰《深入理解 AI Agent：设计原理与工程实践》，v2.0（2026-09-19），第三章“用户记忆和知识库”。

阅读范围：书内第 72-101 页，对应 PDF 文件第 80-109 页。下文页码均指书内页码。

本文按学习与构建顺序重排原文要点。最后的练习顺序是整理建议；教材示例中的模型和项目名不代表已验证的当前安装配置。
\end{quote}

\section{先抓住本章的主线}

\textbf{先把资料变成有结构、有上下文、可追溯的知识，再让 Agent 把检索当作工具，按问题需要反复查询和验证。}

本章讨论两种持久化知识：用户记忆保存某个用户的偏好、经历和关系；共享知识库保存业务规则、技术文档等领域知识。两者共用分块、检索、压缩等技术，也都需要处理冲突、过期和权限问题。

\begingroup\small
\setlength{\tableusable}{\dimexpr\linewidth-4\tabcolsep\relax}
\begin{longtable}{@{}>{\raggedright\arraybackslash}p{0.2\tableusable}>{\raggedright\arraybackslash}p{0.36\tableusable}>{\raggedright\arraybackslash}p{0.44\tableusable}@{}}
\toprule
\textbf{内容} & \textbf{书中主要解决的问题} & \textbf{对应小节} \\
\midrule\endfirsthead
\toprule
\textbf{内容} & \textbf{书中主要解决的问题} & \textbf{对应小节} \\
\midrule\endhead
\bottomrule\endfoot
用户记忆 & 跨会话记住什么、怎样存、怎样找回 & 3.1 \\
基础 RAG & 如何把文档变成可检索的片段，并据此回答 & 3.2 \\
知识组织与更新 & 如何保留层次、关系、规则边界，以及持续维护 & 3.3.1-3.3.3 \\
智能体化 RAG & 如何让 Agent 自主决定查询和后续检索 & 3.3.4 \\
上下文感知检索 & 如何补回分块丢失的背景信息 & 3.3.5 \\
知识发现 & 如何从大量案例中归纳决策因素和模式 & 3.3.6 \\
\end{longtable}\endgroup

这些方法处理不同环节，可以组合使用。它们并不是必须逐级升级、全部用上的固定阶梯。

\section{基础知识库怎样构建}

\subsection{区分建库阶段与问答阶段}

依据 3.2、3.3.5，可将流程整理为：

\begin{quote}\small
\textbf{建库阶段：}原始文档 \(\rightarrow\) 分块 \(\rightarrow\) 按需补上下文前缀 \(\rightarrow\) 建立向量索引和关键词索引


\textbf{问答阶段：}用户问题 \(\rightarrow\) 两路检索 \(\rightarrow\) 结果融合 \(\rightarrow\) 重排序 \(\rightarrow\) 相关内容进入上下文 \(\rightarrow\) LLM 生成回答


\end{quote}

RAG 的基本组成是检索器和生成器：检索器找出相关资料，生成器根据这些资料回答。外部知识可以独立更新，不需要为每次知识变化重新训练大模型。（3.2，第 80-87 页）

\subsection{第一步：把长文档切成可独立检索的片段}

整篇文档只生成一个向量，会混合多个主题；检索整篇文档也容易带入大量无关内容。书中介绍三种分块方式：

\begingroup\small
\setlength{\tableusable}{\dimexpr\linewidth-4\tabcolsep\relax}
\begin{longtable}{@{}>{\raggedright\arraybackslash}p{0.2\tableusable}>{\raggedright\arraybackslash}p{0.36\tableusable}>{\raggedright\arraybackslash}p{0.44\tableusable}@{}}
\toprule
\textbf{策略} & \textbf{做法} & \textbf{主要取舍} \\
\midrule\endfirsthead
\toprule
\textbf{策略} & \textbf{做法} & \textbf{主要取舍} \\
\midrule\endhead
\bottomrule\endfoot
固定大小 & 按 token 数切分，并保留相邻重叠 & 简单，但可能切断段落、代码或表格 \\
递归／结构感知 & 优先按章节、段落切分，超长时再按更小边界切 & 利用文档结构，是书中介绍的常见生产默认选择 \\
语义切分 & 在相邻句子语义相似度明显下降处切分 & 块内主题更一致，需要额外嵌入计算 \\
\end{longtable}\endgroup

书中给出的调参起点是：\textbf{每块 256-1024 token，相邻块重叠 10\%-20\%}，再通过检索质量实测调整。它不是通用最佳参数：块太小会缺少背景，块太大则会混入多个主题。（3.2.1，第 81 页）

\subsection{第二步：同时建立语义检索和关键词检索}

\textbf{稠密检索}把文本转成向量，用相似度寻找语义接近的内容。它适合同义表达、多语言等场景，但可能漏掉必须精确匹配的代码、编号和名称。

书中以 BGE-M3 等模型说明嵌入，并比较两种近似最近邻索引：ANNOY 适合不常变化的静态数据；HNSW 支持增量更新，适合动态场景，但内存占用更高。嵌入模型与索引算法都影响质量、性能和维护成本。（3.2.2，第 82-83 页）

\textbf{稀疏检索}以关键词匹配为核心。书中的 BM25 教学实现依次进行分词、去停用词、建立倒排索引、计算词项权重。BM25 综合词的稀有程度、词频饱和与文档长度，适合精确名称和技术代码，但难以识别字面不同的同义表达。（3.2.3，第 84-85 页）

\subsection{第三步：融合候选，再精细排序}

书中推荐的混合检索流水线包含三个环节：

\begin{enumerate}
\item \textbf{并行召回}：稠密检索与 BM25 分别返回候选。
\item \textbf{结果融合}：将两路结果合成统一候选池。两路原始分数尺度不同，不能直接相加；书中介绍 RRF，按排名的平滑倒数求和。
\item \textbf{神经重排序}：跨编码器同时读取“查询＋候选文档”，对候选池进一步精排。
\end{enumerate}

RRF 可写成以下公式：
\[\operatorname{RRF}(d)=\sum_{r\in\mathcal{R}}\frac{1}{c+\operatorname{rank}_{r}(d)}.\]
这里用 \code{c} 表示书中常取 60 的平滑常数，避免与返回数量 \code{top-\allowbreak k} 混淆。书中“对前 50 个候选重排序”也是示例，需要按成本和质量调整。

\textbf{融合负责形成候选池，重排序负责提高候选池中的排序质量。} 重排序不替代融合，也只能处理已经召回的内容。（3.2.4，第 86-87 页）

\subsection{第四步：补回分块的上下文}

即使切分合理，孤立片段仍可能缺少公司名、时间和章节背景，例如“该公司第二季度收入增长了 3\%”。

书中的上下文感知检索，在\textbf{索引阶段}用 LLM 为片段生成简短背景前缀，然后将“前缀＋原文”一起用于向量和 BM25 索引。例如补充“ACME 公司 2025 年 Q2 财报·关键业绩指标”，让片段的主体和时间明确。

这同时改善关键词匹配与向量语义表示，代价是额外的建库调用。生成的背景仍需有原始资料支持，不能把推测写成事实。

三个概念要分清：

\begin{itemize}
\item \textbf{上下文感知检索}：建库时为片段补背景，提高可检索性。
\item \textbf{智能体化 RAG}：运行时由 Agent 决定怎样查、是否继续查。
\item \textbf{上下文感知压缩}：第二章的方法，运行时裁剪会话内容以节省窗口。
\end{itemize}

前两者可以结合使用。（3.3.5，第 96-98 页）

\section{知识库怎样组织，才能超越零散文本}

\subsection{原始案例入库之前，要提炼可用知识}

书中用两个例子说明，单纯存入原始案例还不够：

\begin{itemize}
\item \textbf{全局统计}：只召回部分猫的案例，不能可靠计算全库黑猫与白猫的比例；应预先生成准确的统计摘要并索引。
\item \textbf{规则边界}：零散客服工单只记录个案结果，可能缺少“适用于谁、哪些人不适用”的完整规则；应结合全库证据提炼带适用边界的规则卡。
\end{itemize}

关键工作发生在索引之前：归纳、抽象、保留条件与边界。增加 \code{top-\allowbreak k} 无法补出原文中从未明确表达的规则。（3.3.1，第 88 页）

\subsection{按问题类型选择组织方法}

\begingroup\small
\setlength{\tableusable}{\dimexpr\linewidth-4\tabcolsep\relax}
\begin{longtable}{@{}>{\raggedright\arraybackslash}p{0.2\tableusable}>{\raggedright\arraybackslash}p{0.36\tableusable}>{\raggedright\arraybackslash}p{0.44\tableusable}@{}}
\toprule
\textbf{方法} & \textbf{构建方式} & \textbf{擅长的问题与限制} \\
\midrule\endfirsthead
\toprule
\textbf{方法} & \textbf{构建方式} & \textbf{擅长的问题与限制} \\
\midrule\endhead
\bottomrule\endfoot
基础混合检索 & 对片段建立稠密／稀疏索引，并重排序 & 寻找某项信息的具体片段，能覆盖多数基础需求 \\
RAPTOR & 片段作为叶子；语义聚类；LLM 生成父层摘要；递归形成知识树 & 从整体概念逐步深入细节、跨层综合 \\
GraphRAG & 提取实体和关系；构建图谱；发现社区并生成摘要 & 多跳关系、实体消歧；三元组可能丢失条件和时间语义 \\
文件系统范式 & 用目录、条目、摘要、索引和交叉链接组织知识 & 便于阅读、编辑、版本管理，以及 Agent 按路径导航 \\
\end{longtable}\endgroup

RAPTOR 与 GraphRAG 会增加构建和查询成本；书中强调，只有经常需要跨文档综合、多层导航或关系推理时，才值得投入。完整自然语言应保留，结构化元数据和图谱可作为补充索引。（3.3.1，第 89-90 页）

\subsection{Markdown 知识库必须有摘要、入口和横向链接}

书中以 OpenViking 说明按需加载的三层表示：

\begin{itemize}
\item \textbf{L0 摘要}：约 100 token，快速判断相关性。
\item \textbf{L1 概览}：约 2000 token，支持规划与决策。
\item \textbf{L2 全文}：需要深入时再读取原始内容。
\end{itemize}

层次摘要解决纵向导航；条目之间的交叉链接、入口页和索引页解决横向导航。只是把 Markdown 文件堆进目录，仍会形成知识孤岛。

书中要求写入新条目时，先查找并链接相关已有条目，同时更新目录索引，形成双向可达的引用网络。纯文本知识还可以通过 Git 审查、追溯和回滚。（3.3.2，第 91 页）

\section{智能体化 RAG 怎样构建}

\subsection{核心变化：检索由 Agent 主动调用}

\begingroup\small
\setlength{\tableusable}{\dimexpr\linewidth-4\tabcolsep\relax}
\begin{longtable}{@{}>{\raggedright\arraybackslash}p{0.2\tableusable}>{\raggedright\arraybackslash}p{0.36\tableusable}>{\raggedright\arraybackslash}p{0.44\tableusable}@{}}
\toprule
\textbf{对比点} & \textbf{传统 RAG} & \textbf{智能体化 RAG} \\
\midrule\endfirsthead
\toprule
\textbf{对比点} & \textbf{传统 RAG} & \textbf{智能体化 RAG} \\
\midrule\endhead
\bottomrule\endfoot
检索触发 & 固定的问答前置步骤 & Agent 按信息需求调用工具 \\
查询组织 & 通常直接使用用户问题 & 分解子问题、改写关键词、追查关系 \\
结果处理 & 一次检索后生成回答 & 观察结果，判断信息缺口，必要时继续检索 \\
成本与用途 & 流程简洁，适合明确单一的问题 & 更多调用和延迟，适合复杂、多跳问题 \\
\end{longtable}\endgroup

\textbf{智能体化的关键是 Agent 拥有检索决策权，不是把搜索次数固定增加到两次或三次。} 简单问题仍可以一次查完。（3.3.4，第 93-95 页）

\subsection{书中的系统组成与循环}

书中的图 3-13 将系统分为 Agent、工具层和可切换的知识库后端。构建顺序可整理为：

\begin{enumerate}
\item \textbf{先建可用的检索后端}：混合检索、结构化索引或上下文感知检索，都可以作为后端。
\item \textbf{将检索封装成工具}：书中的名称是 \code{knowledge\_\allowbreak base\_\allowbreak search}。同一 Agent 还可以调用 \code{web\_\allowbreak search}、\code{code\_\allowbreak interpreter} 等辅助工具。
\item \textbf{接入 ReAct 循环}：分析需求，调用工具，读取返回结果，再评估信息是否充分。
\item \textbf{根据信息缺口调整查询}：需要时分解子问题、搜索关联条款、核对条件或冲突记录。
\item \textbf{证据充分后综合回答}：将多轮获取的信息组织成有依据的结论。
\end{enumerate}


\begin{figure}[htbp]
\centering
\begin{tikzpicture}[
  node distance=7mm,
  box/.style={draw=accent,rounded corners=2pt,fill=light,text width=5.3cm,minimum height=9mm,align=center,font=\small},
  decision/.style={diamond,draw=accent,fill=light,aspect=3,inner sep=2pt,align=center,font=\small},
  arrow/.style={-{Stealth[length=2.2mm]},draw=accent,thick}
]
\node[box] (q) {用户问题};
\node[box,below=of q] (p) {分析需求与信息缺口};
\node[box,below=of p] (t) {调用检索或辅助工具};
\node[box,below=of t] (o) {读取结果并评估证据};
\node[decision,below=of o] (c) {信息是否充分};
\node[box,below=10mm of c] (a) {综合证据生成回答};
\draw[arrow] (q) -- (p);
\draw[arrow] (p) -- (t);
\draw[arrow] (t) -- (o);
\draw[arrow] (o) -- (c);
\draw[arrow] (c) -- node[right,font=\small] {充分} (a);
\draw[arrow] (c.west) -- ++(-2cm,0) |- node[pos=0.25,left,font=\small,align=right] {不足\\调整查询} (p.west);
\end{tikzpicture}

\caption{智能体化 RAG 的检索与验证循环（依据教材图 3-12、3-13 整理）}
\end{figure}


实验 3-8 的复杂司法问答示例，先检索各子问题，再追查它们之间的关系；首次命中若干相关片段后，Agent 仍会检查是否遗漏关键联系。教材借此说明迭代检索的价值，并未要求所有问题都走复杂流程。（第 94-95 页）

\subsection{检索资料必须保持“数据”的身份}

外部文档可能包含提示注入或被投毒的知识。书中提出两层边界：检索结果要标注来源、明确它是参考资料；转账、删除、发信等动作要经过独立授权判断，不能只因文档中的文字而执行。（3.3.4，第 94 页）

\section{知识库怎样持续更新}

\subsection{分开原始证据、审核知识和服务索引}

\begingroup\small
\setlength{\tableusable}{\dimexpr\linewidth-4\tabcolsep\relax}
\begin{longtable}{@{}>{\raggedright\arraybackslash}p{0.2\tableusable}>{\raggedright\arraybackslash}p{0.36\tableusable}>{\raggedright\arraybackslash}p{0.44\tableusable}@{}}
\toprule
\textbf{层次} & \textbf{保存的内容} & \textbf{作用} \\
\midrule\endfirsthead
\toprule
\textbf{层次} & \textbf{保存的内容} & \textbf{作用} \\
\midrule\endhead
\bottomrule\endfoot
原始证据层 & 对话、运行轨迹、原始文档 & 保留原始来源，支持核查 \\
知识层 & 提炼后的 Markdown、规则或代码 & 可修订、可审核、可版本管理 \\
服务层 & 分块、摘要、向量及其他检索索引 & 从已审核版本生成，可重建 \\
\end{longtable}\endgroup

\textbf{已审核的知识是来源，检索索引是派生物。}（3.3.3，第 92 页）

\subsection{增量更新：把知识变更当作 PR}

新证据到来后，书中采用“提议者—审核者”流程：

\begin{enumerate}
\item Proposer Agent 搜索已有知识，在工作分支提出最小、完整的修改，同时维护链接、时间和证据引用。
\item Reviewer Agent 独立对照旧知识、差异和原始证据，检查断言、限定条件、冲突及过度改写。
\item 双方按具体反馈迭代；设置迭代或成本上限，未收敛时转人工审核。
\item 批准并通过格式、链接、元数据、权限等检查后合入，再重建受影响的索引。
\end{enumerate}

书中特别强调：这两个角色应有主动搜索、追溯和检查能力，不只是两次固定的摘要调用。优先使用能力相近、来自不同模型家族的 Agent，并限制各角色的写入权限。审核者应能在授权范围内查完整证据，不能只看提议者选出的片段。（第 92 页）

\subsection{定期整理：回到原始数据重审全库}

周期整理包括去重、去旧、合并、调整目录和链接，以及核对摘要是否遗漏事实、否定词或时间条件。大型库可分批处理，但需要覆盖清单。

冲突不能简单按“最新一条”覆盖：先确认是否在不同时间、对象或条件下分别成立；证据不足时保留冲突和待确认状态。全量整理也要经过审核，再重建索引、回放典型检索与问答用例。

检索时还要利用版本、生效／失效时间过滤旧知识，并按调用者权限、租户范围过滤资料，避免越权内容进入上下文。（第 93 页）

\section{怎样把同一套技术用于用户记忆}

用户长期记忆是从对话中提炼的稳定信息；原始对话轨迹是可追溯的事件记录。两者用途不同。记忆提取遵循选择性、抽象化和结构化。（3.1，第 72-75 页）

书中的四种格式，从简单到复杂为：

\begingroup\small
\setlength{\tableusable}{\dimexpr\linewidth-2\tabcolsep\relax}
\begin{longtable}{@{}>{\raggedright\arraybackslash}p{0.3\tableusable}>{\raggedright\arraybackslash}p{0.7\tableusable}@{}}
\toprule
\textbf{格式} & \textbf{特点} \\
\midrule\endfirsthead
\toprule
\textbf{格式} & \textbf{特点} \\
\midrule\endhead
\bottomrule\endfoot
Simple Notes & 单条原子事实，便宜，但关联容易丢失 \\
Enhanced Notes & 完整叙事段落，语义完整，但冗余、更新复杂 \\
JSON Cards & 分类字段，支持局部更新，但分类可能过于刚性 \\
Advanced JSON Cards & 加入背景、主体、关系和时间，便于消歧，维护成本较高 \\
\end{longtable}\endgroup

实验 3-9 将完整对话历史作为知识库，给 Agent 配置 \code{search\_\allowbreak user\_\allowbreak memory} 工具。用户有两辆车时，Agent 可沿首次结果中的线索继续检索另一辆车，而不是随意猜测。但孤立对话块仍可能难以表达连续修改和时间关系。（第 95-96 页）

本章最终汇合为\textbf{双层记忆架构}：少量关键事实用 Advanced JSON Cards 常驻上下文，提供全局概览；大量原始对话用上下文感知检索按需召回，提供精确细节。

两层配合，既有机会发现跨会话关联，也能回到原始记录验证。书中据此划分三个验收层次：基础回忆、多会话检索、主动服务。（第 73-74、97-98 页）

另外，书中讨论了可执行记忆：把状态写成带类型对象、把聚合和约束写成函数，减少让 LLM 做精确统计或规则检查的负担。它是另一种知识表示选择，不是构建基础 RAG 的必选项。（3.1.4，第 76-77 页）

\section{怎样判断构建效果}

书中要求在带标注的测试查询集上衡量检索质量，并通过组件对比观察各环节的贡献。（3.2.4，第 87 页；实验 3-8，第 95 页）

\begingroup\small
\setlength{\tableusable}{\dimexpr\linewidth-2\tabcolsep\relax}
\begin{longtable}{@{}>{\raggedright\arraybackslash}p{0.3\tableusable}>{\raggedright\arraybackslash}p{0.7\tableusable}@{}}
\toprule
\textbf{指标／对比} & \textbf{关注点} \\
\midrule\endfirsthead
\toprule
\textbf{指标／对比} & \textbf{关注点} \\
\midrule\endhead
\bottomrule\endfoot
书中的 \code{recall@k} & 前 k 个结果是否至少命中一篇相关文档 \\
MRR & 第一个相关结果是否排得足够靠前 \\
nDCG & 整个结果列表的相关程度和排序质量 \\
传统分块 vs 上下文感知分块 & 补背景后是否更准确命中主体、时间与意图 \\
传统 RAG vs 智能体化 RAG & 简单／复杂问题上的质量与响应成本差异 \\
\end{longtable}\endgroup

\textbf{指标口径要注意}：第 87 页脚注明确，书中 \code{recall@k} 实际采用命中率（hit rate／success@k）；标准召回率是“召回的相关文档数 \(\div\) 全部相关文档数”。多篇相关文档场景下，两者不同。

书中引用的效果与成本属于相应研究或实验条件，不应直接作为自己系统的效果保证。尤其需要检查多跳问题、冲突记录、失效知识和权限过滤，而不只看几个成功回答。

\section{建议的学习与练习顺序}

以下根据本章内容整理，不是书中逐字给出的部署步骤：

\begin{enumerate}
\item \textbf{先做基础问答库}：结构感知分块，接入稠密检索、BM25、融合和重排序，建立可核对的查询集。
\item \textbf{补足片段背景}：加入上下文前缀，用同一查询集比较前后差异。
\item \textbf{改善知识组织}：增加统计摘要、规则卡、入口和链接；确有层次或关系需求时再加入 RAPTOR／GraphRAG。
\item \textbf{把搜索交给 Agent}：封装检索工具，接入 ReAct，让它根据缺失信息迭代查询。
\item \textbf{加入维护流程}：保存证据、版本和权限信息，让更新经过审核，再发布派生索引。
\item \textbf{扩展到用户记忆}：结合常驻事实概览与按需对话检索，按三个记忆层次评估。
\end{enumerate}

教材相关实验：3-4 \code{dense-\allowbreak embedding}、3-5 \code{sparse-\allowbreak embedding}、3-6 \code{retrieval-\allowbreak pipeline}、3-7 \code{structured-\allowbreak index}、3-8 \code{agentic-\allowbreak rag}、3-9 \code{agentic-\allowbreak rag-\allowbreak for-\allowbreak user-\allowbreak memory}、3-10 \code{contextual-\allowbreak retrieval}、3-11 上下文感知用户记忆。

本章还有两项进阶内容：从案例数据中提取标准化 JSON，进行因子分析与聚类，形成可供 Agent 使用的决策模式（实验 3-12，第 98-99 页）；以及多模态记忆的探索（第 100 页）。前者强调知识发现，后者涉及原始媒体、嵌入或特定模型结构，不宜当作基础文本 RAG 的必备步骤。
\end{document}
